\section{Repair of the full-completion inputs to E.42}
\label{sec:completion-repair}

\begin{result}{A.6}{Proved}{EXIT\_TOWER E.42--E.43, with the missing completion arguments supplied}
Equation (E.200) is precisely Lochak's compatible-pair equation (5.1).
The equality of its two parameters follows from the genus-one calculation
of A.1 and the parameter detection proved in Section~\ref{sec:tower-audit}.
The completed exit system used in E.42 is the completion by all finite
categorical congruences, and its endpoint groups are the full profinite
mapping-class groups, in every genus.  Restriction to its labelled
genus-zero system and the extension of E.39 are inverse on coherent
operation-preserving equivalence classes.  The absent proofs of E.29,
E.32, E.33, E.35 and E.38 are replaced by the arguments below; in
particular, no congruence-subgroup-property assertion in genus at
least three is required.
\end{result}

The missing source files for E.29, E.32, E.33, E.35 and E.38 are an
omission in the supplied proof archive.  We supply the particular
completion arguments needed by E.39--E.42 here.  Throughout, completion
means the inverse limit of all finite categorical congruences on a
finite orbit skeleton.  All mapping-class groups are completed with
respect to all their finite quotients.

\subsection{Finite convex extensions and exact fibre products}

Let $\mathcal D$ be a full convex subcategory of a finite-object
exit-category skeleton $\mathcal C$: every composable path with both
endpoints in $\mathcal D$ has all its intermediate objects there.
Every finite quotient of $\mathcal D$ extends to a finite quotient
of $\mathcal C$.  Keep that quotient on the Hom-sets of $\mathcal D$;
replace every nonempty Hom-set having an endpoint outside
$\mathcal D$ by a singleton, and retain the empty Hom-sets.
Composition inside $\mathcal D$ is its prescribed composition.
Every other composite has its forced singleton value.  A composite
with both endpoints inside cannot pass through an outside object,
by convexity.  Hence no outside composition imposes an additional
relation inside.  This verifies identities and associativity and
proves the extension assertion.  The same argument applies to a finite
union of convex components: a path leaving that union and returning
would violate its stated convexity.  Completion consequently restricts
to the actual full completion on a convex subcategory.

We also record the exact fibre-product argument used below.  For a
surjection of discrete groups $H\to G$ with kernel $N$, and a
homomorphism $K\to G$, put $P=H\times_G K$.
In each application here, $N$ has its full induced profinite topology
in $H$.  Projection $P\to H$ then gives the full induced topology
on the copy of $N$ in $P$ as well.  Right exactness of profinite
completion gives exact sequences
\[
\begin{split}
1&\longrightarrow\widehat N\longrightarrow\widehat H
                      \longrightarrow\widehat G\longrightarrow1,\\
1&\longrightarrow\widehat N\longrightarrow\widehat P
                      \longrightarrow\widehat K\longrightarrow1.
\end{split}                                                     \tag{TC.1}
\]
The initial injections follow from the full induced topologies.
The kernels of the right arrows are the closures of their discrete
kernels, by right exactness, so no separability assertion has been
substituted for this argument.  The profinite fibre product
$\widehat H\times_{\widehat G}\widehat K$ has the second exact
sequence in (TC.1) too.  The natural homomorphism from $\widehat P$
induces the identity on its kernel and quotient; it is therefore an
isomorphism.  For a finite-index image in place of a surjection, replace
$G$ by that image and $K$ by its inverse image.  The latter is a
finite-index subgroup, whose completion is the corresponding open
subgroup.  The same proof then applies, retaining the finite quotient
of components.  This is the exact form needed in the coarse flag
calculations of E.40.

\subsection{Crossings and adjacent flags}

Let $\Gamma$ be a mapping-class group and $\sigma$ an actual
multicurve.  Write $C=\operatorname{Stab}_\Gamma(\sigma)$,
$\Delta=\langle\tau_c:c\in\sigma\rangle$ and $H=C/\Delta$.
The discrete crossing to the open object is the transitive left
$\Gamma$-set $\Gamma/\Delta$.  Every finite categorical congruence
restricts on this set to a finite equivariant quotient, and its action
factors through a finite quotient of $\Gamma$.  Conversely, every
finite quotient $\Gamma/N$ induces a quotient of the whole
transporter category: replace each representative $g$ by its image
and each source multitwist group by its image.  This is well-defined
under composition because
$g\sigma_i\supseteq\sigma_j$ implies
$g^{-1}\Delta_{\sigma_j}g\subseteq\Delta_{\sigma_i}$.
Thus the crossing has exactly the topology of the inverse system
$\Gamma/(N\Delta)$, and
\[
       \widehat{\operatorname{Hom}}(\sigma,\varnothing)
                  =\widehat\Gamma/\overline\Delta.         \tag{TC.2}
\]
Every endpoint group has its full completion: extend any finite quotient
of that one-object convex subcategory by the preceding construction.

Boggi's Theorem~4.5 \cite{ET_Boggi2014}, specifically the injectivity of the multitwist
lattice in its second exact sequence, supplies
$\widehat\Delta\hookrightarrow\check\Gamma$.
The map factors through $\widehat\Gamma\twoheadrightarrow
\check\Gamma$, so its injection also proves
$\widehat\Delta\hookrightarrow\widehat\Gamma$.
This is the precise direction in which the cited procongruence result
is used.  In particular $\overline\Delta=\widehat\Delta$ in (TC.2).
The finite-quotient diagonal argument of the preceding section gives
\[
        \widehat C\hookrightarrow\widehat\Gamma\times\widehat H.
\]
Lifting the $\widehat H$ coordinate and then the residual
$\widehat\Delta$ coordinate computes the whole stabilizer of the
pinned crossing as $\widehat C$.

For an adjacent incidence $\sigma\supset\tau$, normalize at the
retained curves $\tau$ and keep their branch pinnings.  This cuts the
target into its actual components.  The incidence is a crossing to
the open object of one component; the other components contribute
identity factors.  Its two endpoint groups and incidence group are
\[
 G_0=\operatorname{Stab}(\sigma)/\Delta_\sigma,
 \quad G_1=\operatorname{Stab}(\tau)/\Delta_\tau,
 \quad K=\operatorname{Stab}(\sigma,\tau)/\Delta_\tau.
\]
The map $K\to G_0$ has finite-index image, since it merely retains
a specified subset of the finite curve set of $\sigma$; its kernel is
$\Delta_\sigma/\Delta_\tau$.  This kernel is the full multitwist
lattice in the normalized target components, and has its full induced
topology there by the same injectivity just used.  The diagonal argument,
with the finite-index image retained, gives
\[
 \widehat K\hookrightarrow\widehat G_1\times\widehat G_0,
 \qquad
 \widehat{\operatorname{Hom}}_{\rm incidence}
     = (\widehat G_1\times\widehat G_0)/\widehat K.          \tag{TC.3}
\]
At the level of a two-object adjacent interval, every finite equivariant
quotient of its crossing biset is a finite categorical quotient;
there is no composition of two noninvertible arrows in this interval.
Convex extension therefore proves that (TC.3) has the topology induced
by the full exit category.

Here is the precise finite-pinning comparison needed to use
normalization inside that category.  On its finite skeleton let $P(i)$
be the finite set of retained submulticurves of the specified type,
with their branch pinnings.  An arrow represented by $g:i\to j$
pulls a target pin $\rho$ back to $g^{-1}\rho$; source twists have no
effect on this pullback.  Thus $P$ is a finite contravariant functor.
An arrow has a unique lift to the category of elements of $P$ for each
target pin.  For a finite quotient $F$ of this marked category, record
on a downstairs arrow $g$ its finite map $P(g)$ and the finite array
$(F(g_\rho))_{\rho\in P(j)}$.  For a composite $hg$, the entries
are $F(h_\rho)F(g_{P(h)\rho})$.  Equality of these arrays is a finite
categorical congruence downstairs, and its pullback refines $F$.
The converse pullback is immediate after retaining $P$.
This proves cofinality and hence commutation of completion with this
particular finite marking construction.  Cutting and extending over
the retained collars are inverse functors on its component diagrams;
their collar twists are killed on both sides.  An intermediate
contraction cannot erase a node retained by the final target, so the
component interval is convex.  All finite quotients of the normalized
component diagrams are therefore retained in the original completion.

Apply (TC.1) successively to the incidence stabilizers of a specified
adjacent flag.  Every kernel that occurs is the residual multitwist
lattice, with the full topology proved above; every image imposing
an order or pin on a finite curve set has finite index.  The stabilizer
of the completed flag is consequently the full completion of its
discrete stabilizer, with precisely the original quotient maps to its
endpoints and arrows.  There are only finitely many matching incidence
orbits for a fixed finite object list.  Their completed group orbits
are compact and closed in the product of Hom-spaces.  Their finite union
contains the product of the dense discrete Hom-spaces, hence is the
whole product.  This proves the required flag assertion without
postulating density in a square equalizer.  It repairs E.32.

\subsection{Pair-tail normalizers in every genus}

We first prove the peripheral fact used in the repair.  For a hyperbolic
punctured surface $S$, let $P=\widehat{\pi_1(S)}$ and let $I$ be a
peripheral procyclic subgroup.  Then
\[
                         N_P(I)=I.                          \tag{TC.4}
\]
Indeed, take an element $h\notin I$ normalizing $I$.  There is an open
subgroup $U\leq P$ containing $I$ and excluding $h$: separate the
compact subgroup $I$ from $h$ in a finite quotient and take the inverse
image of the image of $I$.  The corresponding finite surface cover
has at least three punctures when the base has genus zero.  In positive
genus its compactification has positive genus by Riemann--Hurwitz.
In that case pass further to the kernel of the map to the mod-two
homology of the compactification.  This is an unramified compact cover
of degree at least four, contains all cusp inertia, and still excludes
$h$.  Denote the resulting subgroup by $V$; its surface has at least
three punctures.

The cosets $V$ and $Vh$ determine distinct unramified cusps over the
chosen base cusp.  Their inertia groups are conjugate in $V$ to the
same subgroup $I$.  To interpret a profinite coset as a topological
cusp, choose a discrete representative $r\in Vh$; writing $r=vh$
gives $rIr^{-1}=vIv^{-1}$.  On the other hand, in the homology of a
surface with at least three punctures the peripheral classes have
exactly their single sum relation.  The procyclic lines generated by
two distinct peripheral classes are distinct.  Conjugate inertia
subgroups would give the same line, a contradiction.  This proves
(TC.4).  It also proves that a primitive generator of a nonabelian
free profinite group has its own procyclic group as centralizer.
Two different free basis elements consequently show that such a free
profinite group has trivial centre.

The closed case also has the required centrelessness.  Let
$\Pi=\pi_1(\Sigma_g)$ with $g\geq2$, and let $z$ be central in
$\widehat\Pi$.  For an arbitrary finite quotient $q:\Pi\twoheadrightarrow Q$,
put $U=\ker q$ and choose a prime $p$ not dividing
$|Q|(2g-2)$.  The finite covering surface has genus
$g_U=1+|Q|(g-1)$.  Its $Q$-equivariant cellular complex, obtained by
lifting the usual one-vertex, $2g$-edge, one-face cell structure, has
regular chain modules of ranks $1,2g,1$.  Both $H_0$ and $H_2$ are
trivial $Q$-modules, since the deck transformations preserve orientation.
Consequently the trace of a nonidentity $h\in Q$ on
$H_1(U,\mathbb F_p)$ is $2$, whereas the trace of the identity is
$2+|Q|(2g-2)$; these are different in $\mathbb F_p$.  Thus the deck
action of $Q$ on this homology is faithful.  Finite-index subgroups
have their full induced profinite topology.  Since $z$ centralizes
$\widehat U$, it acts trivially on this finite homology quotient.
The action factors through $Q$, so $q(z)=1$.  This holds for every
finite quotient $q$, proving $Z(\widehat\Pi)=1$.

For every stable surface of type $(g,n)$, including $n=0$ and $g\geq2$,
the discrete Birman sequence for an additional labelled point is
\[
1\longrightarrow\pi_{g,n}\longrightarrow\Gamma_{g,n+1}
                     \longrightarrow\Gamma_{g,n}\longrightarrow1.
\]
Its kernel has the full induced topology.  Namely, use the extra point
as basepoint to obtain
$\Gamma_{g,n+1}\to\operatorname{Aut}(\widehat\pi_{g,n})$.
Characteristic finite quotients of the finitely generated surface
group (free when $n>0$) give continuity from the full profinite completion.  On
point pushes this is the map to inner automorphisms.  The centre of
$\widehat\pi_{g,n}$ is trivial by the two preceding arguments, so that
inner-automorphism map is injective.  Its factorization through
$\widehat\Gamma_{g,n+1}$ proves the asserted injection.  Right
exactness now proves the full completed Birman sequence.

Attach one pair-tail at an old labelled puncture $p$, with new point
$x$, and let $C$ stabilize its separating curve $t$ in
$\Gamma_{g,n+1}$.  Cutting off that two-punctured disc identifies $C$
with the group of the core with $p$ blown up to one fixed boundary.
Forgetting $x$ restricts to the marked-capping surjection
\[
             C\longrightarrow\Gamma_{g,n},\qquad
                    \ker=\mathbb Z\tau_t.
\]
The twist is the actual peripheral element of the Birman kernel.
The full diagonal topology, together with the forgetting map on the
ambient group, proves $\widehat C\hookrightarrow
\widehat\Gamma_{g,n+1}$: its second diagonal coordinate is already
a function of its first.  Moreover
$\widehat C\cap\widehat\pi_{g,n}=\widehat{\mathbb Z}\tau_t$.
A normalizer of this twist group can be multiplied by an element of
$\widehat C$ so that its forgetting image is one.  It then lies in the
Birman kernel and (TC.4) puts it in the twist group itself.  Since all
these mapping classes preserve orientation and the two labels, $C$
centralizes the twist.  Thus
\[
 N_{\widehat\Gamma_{g,n+1}}(\widehat{\mathbb Z}\tau_t)
 =C_{\widehat\Gamma_{g,n+1}}(\widehat{\mathbb Z}\tau_t)
 =\widehat C.                                                \tag{TC.5}
\]

Now attach $r$ tails at distinct old punctures, with new points
$x_1,\ldots,x_r$, curves $t_1,\ldots,t_r$ and twist lattice $\Delta$.
The joint stabilizer has its full topology, by induction on $r$.
In fact the single-$t_r$ stabilizer maps by marked capping onto the
group with $x_r$ forgotten, with kernel $\mathbb Z\tau_{t_r}$.
The joint stabilizer is the inverse image of the joint stabilizer of
the other $r-1$ curves.  To check this discrete equality, isotope a
representative fixing those curves and the retained point $p_r$;
a small disc at $p_r$ can be chosen disjoint from the curves, so it
lifts to a boundary-fixing representative.  All lifts differ by the
boundary twist, which fixes the other curves.  The exact fibre-product
calculation (TC.1) and the induction injection prove the same
inverse-image equality after full completion.  This also identifies
the intersection of the completed single-tail stabilizers with the
completion of their discrete intersection.

An element normalizing $\widehat\Delta$ acts on it by a matrix in
$\operatorname{GL}_r(\widehat{\mathbb Z})$.  Forget every $x_j$
except $x_i$.  The induced map sends $\tau_{t_i}$ to the single-tail
twist and kills all $\tau_{t_j}$ with $j\ne i$.
Applying this map to the conjugation equation for $\tau_{t_j}$ forces
the matrix entry in row $i$, column $j$ to be zero when $i\ne j$.
Thus the matrix is diagonal and the normalizer preserves every
individual twist line.  Equation (TC.5) and the intersection statement
then give
\[
 N_{\widehat\Gamma}(\widehat\Delta)
 =C_{\widehat\Gamma}(\widehat\Delta)
 =\widehat{\operatorname{Stab}_\Gamma(t_1,\ldots,t_r)}.
                                                               \tag{TC.6}
\]
This proves the full-completion normalizer assertion of E.33.

There is also a direct proof of its Boolean-category assertion.  For
subsets $I\supseteq J$ of the tail set, every transporter from $t_I$
to $t_J$ fixes each $t_j$, $j\in J$: its pair of fresh labels singles
it out among the source curves.  Writing $C_J=\operatorname{Stab}(t_J)$,
we have, with their actual inclusions,
\[
 \operatorname{Hom}(t_I,t_J)=C_J/\Delta_I,
 \qquad
 \widehat{\operatorname{Hom}}(t_I,t_J)
                       =\widehat C_J/\widehat\Delta_I.       \tag{TC.7}
\]
For the topology in (TC.7), the target endpoint group acts transitively,
so every finite categorical quotient factors through a finite quotient
of it.  Conversely the ambient finite transporter quotients induce
all the needed finite quotients of $C_J$, because $C_J$ has its full
topology in $\Gamma$.  This proves equality of the two topologies.
The map from (TC.7) to $\widehat\Gamma/\widehat\Delta_I$ is injective.

An equivalence with identity open-group action sends the pinned crossing
from $t_I$ to the open object to a right translate by a unique element
of $\widehat C_I/\widehat\Delta_I$, by (TC.6).  Use that element as
the comparison at $t_I$.  All pinned crossings and their endpoint
actions are then fixed, since the right action is faithful and free.
Composition into the open crossing is injective by (TC.7), so all
remaining arrows are fixed.  The uniqueness at each endpoint makes
these comparisons coherent.  This proves the stated Boolean rigidity,
including every crossing map, without a general multicurve-normalizer
assumption.

\subsection{The exact adjacent presentation}

Before passage to an orbit skeleton, use all actual multicurves as
objects.  Introduce isomorphisms $u_g:\sigma\to g\sigma$ for
$g\in\Gamma$ and one-curve deletion arrows
$d_{\sigma,c}:\sigma\to\sigma\setminus\{c\}$.
Impose the multiplication equations on the $u_g$, their equivariance
with deletions, $u_\delta=1$ at $\sigma$ for $\delta\in\Delta_\sigma$,
and the diamond deleting two distinct curves in either order.
Every word can be moved, by equivariance, into a single $u_g$ followed
by deletions of precisely the curves in $g\sigma\setminus\tau$.
Diamonds make the order of these deletions irrelevant.  The resulting
normal form is exactly
\[
       g\Delta_\sigma,
          \qquad g\sigma\supseteq\tau.                     \tag{TC.8}
\]
Changing $g$ by a source twist leaves the word unchanged by the stated
source relation.  Conversely two equal normal forms differ by such a
twist.  Multiplication of representatives gives the transporter
composition, since the twists at a target lift into the source
twist lattice.  These assignments construct inverse functors between
the presented category and the actual transporter category.

Choose one representative per finite orbit and choose a transporter
from each actual multicurve to its representative.  Conjugating the
above generators by these chosen transports gives the skeletal
endpoint groups, adjacent incidence bisets and pinned diamonds.
The inverse includes the representatives and uses those transports
as its natural isomorphism.  It retains every actual pin and every
incidence orbit.  A finite assignment of the completed endpoint
groups and adjacent bisets respecting these equations restricts to
a functor on the discrete category by (TC.8), and its finite kernel
is a categorical congruence.  Conversely every such congruence gives
that assignment.  Taking the inverse limit over finite targets proves
the universal profinite presentation with the same continuously
transported diamonds.  This supplies E.38 in precisely the finite
quotient sense used by E.39.

\subsection{Framed square equalizers and the sewing section}

The repaired ingredients suffice for the square argument in E.40;
we give the complete reason that no extra completed squares occur.
Attach distinct pair-tail tags $t$ to a bordered core and write
$A_i=t\cup\gamma_i$, $B_i=\gamma_i$, with
$\gamma_0\supseteq\gamma_1$ and fixed vertical tag-smoothing arrows
$\eta_i:A_i\to B_i$.  Set
\[
 C_i=\operatorname{Stab}(t,\gamma_i)/\Delta_{\gamma_i},
 \qquad S=\operatorname{Stab}(t,\gamma_0,\gamma_1)
                                      /\Delta_{\gamma_1}.
\]
Discretely a commutative square is represented by one element
$c\in\operatorname{Stab}(t)$ and has horizontal representatives
$c\Delta_{t\cup\gamma_0}$ and $c\Delta_{\gamma_0}$.
A tag cannot be exchanged with an interior essential core curve,
because its pair of new labels distinguishes it.  The two square
composites agree precisely modulo
$\Delta_t\Delta_{\gamma_0}$, which proves this description and
its inverse, the bottom representative.

At the endpoints, (TC.6) gives the full groups
$\operatorname{Aut}(\eta_i)=\widehat C_i$, with injective target
projection and surjective source projection of kernel
$\widehat\Delta_t$.  For the common composite $z_0:A_0\to B_1$,
normalize at $\gamma_1$ and use the finite-pinning comparison above.
Its orbit is $E/\widehat S$, where
$E=\widehat{\operatorname{Aut}(B_1)}\times
\widehat{\operatorname{Aut}(A_0)}$.
Each of its two indicated intermediate subsets is intrinsic, because
tags cannot be exchanged with core curves.  Discretely its specified
two-step factorization stabilizer projects isomorphically onto $S$;
the missing intermediate coordinate is the same representative modulo
the intermediate twist lattice.  Applying (TC.1) with kernel
$\Delta_{\gamma_0}/\Delta_{\gamma_1}$ on one side and $\Delta_t$
on the other proves the same assertion after completion.

There are finitely many matching factorization orbits, and their compact
completed orbits cover the product of composable Hom-spaces by the
ordinary-product density argument above.  Thus the two factorization
spaces of this common-composite orbit are exactly
\[
 (E\times\widehat{\operatorname{Aut}(B_0)})/
       \operatorname{Graph}(\widehat S\to
                          \widehat{\operatorname{Aut}(B_0)}),
\quad
 (E\times\widehat{\operatorname{Aut}(A_1)})/
       \operatorname{Graph}(\widehat S\to
                          \widehat{\operatorname{Aut}(A_1)}).
                                                               \tag{TC.9}
\]
Composition in both spaces is the projection to $E/\widehat S$.
Write a common composite as $(g_1,h_0)z_0$ and use middle coordinates
$x,y$ in these two spaces.  Requiring the vertical arrows to be the
fixed $\eta_0,\eta_1$ is exactly
\[
                 (x,h_0)\in\widehat C_0,
                  \qquad(g_1,y)\in\widehat C_1.
\]
The resulting square space is therefore
\[
                 (\widehat C_1\times\widehat C_0)/\widehat S.
                                                               \tag{TC.10}
\]
Its stabilizer has full topology by the diagonal argument, since its
source map has finite-index image and residual full multitwist kernel.
The same crossing computation on the bordered core shows that (TC.10)
is its full framed-transporter completion.  The two horizontal
projections are the original quotient maps.  Refinements and
compositions agree on representatives and then, by density in products,
on the completed Hom-spaces.  This proves E.40's actual equalizer
comparison, rather than asserting an abstract interchange between
completion and equalizers.

For the sewing recovery, choose two labelled boundaries $a,b$ of the
core, attach a tripod to each, clutch one external leg of each tripod,
and retain their other legs as distinct companion labels.  Let $i,j$
be the core interfaces and $m$ the middle collar.  The flag is
\[
 \{i,m,j\}\cup\gamma\longrightarrow\{m\}\cup\gamma
                                      \longrightarrow\gamma.
\]
A stable trivalent genus-zero vertex cannot be created by contraction:
for a connected contracted subgraph $H$ the exact identity is
\[
3g(v)-3+\operatorname{val}(v)
 =\sum_{w\in V(H)}(3g(w)-3+\operatorname{val}(w))+|E(H)|.
\]
All summands are nonnegative.  Each companion therefore singles out
its original tripod, and every transporter of these flags preserves
$i,m,j$ individually.  Cutting the two tripods leaves a framed core
transporter $p$ and an independent middle twist $n\in\mathbb Z$.
Each inner tripod twist is identified with its core boundary twist,
each companion twist is killed by its marked cap, and the two middle
leg twists are identified.  These relations give, on all arrows,
\[
                         L=P\times B\mathbb Z_m.            \tag{TC.11}
\]
Here $P$ is the full framed exit category of the core.
The square calculation (TC.10), including the retained finite pins,
proves the completed version of (TC.11).  On a single orbit its middle
coordinate is the difference $n_1-n_0$ of the endpoint coordinates;
the stabilizer has the diagonal middle coordinate.  Explicitly,
\[
 \frac{(\widehat P_1\times\widehat{\mathbb Z})
                  \times(\widehat P_0\times\widehat{\mathbb Z})}
      {\widehat S_P\times\operatorname{diag}(\widehat{\mathbb Z})}
 \longrightarrow
 \bigl((\widehat P_1\times\widehat P_0)/\widehat S_P\bigr)
                          \times\widehat{\mathbb Z},
\quad[(p_1,n_1),(p_0,n_0)]\mapsto([p_1,p_0],n_1-n_0)
\]
is an isomorphism with inverse given by $n_0=0,n_1=n$.
Finite categorical quotients also prove this product completion:
take a common multiple of the orders of the images of the central
integer at the finitely many objects, and take the quotient induced
on $P$.  These product congruences are cofinal.

Both coordinate maps are retained operations.  Normalizing the retained
middle node gives $q(p,n)=p$.  Ordering the middle smoothing first and
normalizing the two interfaces selects the four-point middle crossing,
giving $r(p,n)=n$.  Changing that order uses precisely the pinned
diamonds of (TC.8).  The simultaneous section is uniquely determined
on every Hom-set by $qs=1_P$, $rs=0$, hence is $(p,0)$.
Forgetting the two companions concatenates the three collars; in their
normal coordinates and winding coordinates the maps are
\[
                   q_e=q_iq_mq_j,
                 \qquad(u_i,u_m,u_j)\mapsto u_i+u_m+u_j.
\]
On the simultaneous section the middle winding is zero, so the bottom
arrow is exactly the smooth sewing of the original framed arrow.
For several sewings, all middle coordinates are zero; the same
characterization proves independence of order, associativity and
self-sewing coherence.  Ordinary caps are recovered by forgetting
both labels of the corresponding tail and stabilizing.  The cylinder
uses the single four-point crossing, with its two boundary twists
identified with that crossing's one positive meridian.  This repairs
the groupoid sewing extraction E.35 as well as the whole-arrow version
used in E.42.

\subsection{The whole-arrow kernel and its exact dependencies}

Every adjacent crossing has a local four-holed-sphere or one-holed-torus
support after normalizing the retained nodes and choosing a pants
refinement.  Sew the remaining pants and retain the old nodes by nodal
clutching.  This sends the local pinned crossing to the original
crossing, with its actual endpoint maps.  The one-holed-torus crossing
itself comes from the four-holed sphere: take the two sides of its
node to have boundary labels $\{a,a_0\}$ and $\{b,b_0\}$,
self-sew $a,b$, and cap $b_0$.  Cutting the original node then leaves
a pair of pants with boundaries $a_0,c^+,c^-$, proving that it is the
nonseparating one-holed-torus node.  The plumbing coordinate at that
node is untouched, so its meridian has no extra winding.
Thus four-point crossings, the recovered whole-arrow sewing, caps and
nodal clutching generate every adjacent pinned crossing.

The recovered boundary groupoids are the marked-pants groupoids by
extension over the fixed tripods and restriction to their cores.
Their operadic comparison retains parametrizations and relabellings.
We spell out the strictification, including graph automorphisms.
The object modular operad is the free one on the
$\mathfrak S_3$-set $\mathfrak S_3/C_3$ of the two cyclically ordered
three-boundary corollas.  Concretely, an object is an isomorphism class
of a connected trivalent graph, with labelled external legs and a
cyclic ordering of the three flags at every vertex.  This is exactly
the graph quotient of
\cite[Definitions~2.24 and 3.1]{ET_BonattoRobertson2026}.

Let $H_P:P\longrightarrow F(P)$ be the component at either corolla
of the chosen genus-zero operadic comparison.  Its equivariance says
$H_{\rho P}=\rho H_P$ for $\rho\in\mathfrak S_3$.
In particular, the three elements of $C_3$ fix $H_P$ when they fix
$P$.  For a decorated graph $G$, form $H_G$ by applying all its
gluing and contraction operations simultaneously to the family
$(H_{P_v})_{v\in V(G)}$, together with the already recovered
structural comparisons of $F$.  These structural comparisons
satisfy associativity, interchange and contraction symmetry because
both sides are the unique section with every middle winding zero.

Here is the precise descent through the graph quotient.  A graph
isomorphism $\psi:G\longrightarrow G'$ consists of a permutation
of vertices and, at each vertex, a bijection of its three flags;
the latter is expressed in local labels by
$\rho_v\in\mathfrak S_3$.  Its action on the decorating morphisms
is the permutation of the factors followed by
$H_{P_v}\mapsto \rho_v H_{P_v}=H_{\rho_vP_v}$.
Equivariance of gluing identifies their simultaneous composites.
Changing the representative of a cyclic ordering changes
$\rho_v$ by $C_3$, which fixes the corresponding decorating
morphism.  Reversing an internal edge exchanges the two inputs
of one contraction; contraction symmetry identifies the result.
Permuting the internal edges changes the order of contractions;
interchange identifies the result.  These operations exhaust a
change of labelled graph representative and a change of an
iterated expression for its simultaneous gluing.  Therefore $H_G$
descends to the actual object of the free modular operad, even
when the graph has a nontrivial automorphism group.  It satisfies
\[
 H_{G\circ K}=H_G\circ H_K,\qquad
 H_{\xi_{ij}G}=\xi_{ij}H_G,\qquad
 H_{\rho G}=\rho H_G,
                                                               \tag{TC.12}
\]
with the structural comparisons understood in these equalities.
The cylinder comparison is fixed by its one-dimensional crossing.

For $\phi:G\to G'$, define
\[
                F'(\phi)=H_{G'}^{-1}F(\phi)H_G.
                                                               \tag{TC.13}
\]
The endpoint maps of $H_G$ make $F'$ the identity on objects.
Equations (TC.12), with the recovered structural comparisons,
give strict compatibility of (TC.13) with all operad operations.
Conjugation by fixed morphisms is continuous; it also gives the
operadic natural comparison $F'\Rightarrow F$.
This is the construction of
\cite[Lemma~6.11]{ET_BonattoRobertson2026}, now with the graph
quotient made explicit.

One can check its descent without choosing a cyclic flag at any
vertex.  Replace $\mathfrak S_3/C_3$ by the regular
$\mathfrak S_3$-set, and keep the same Hom-sets after forgetting
the extra flag.  The resulting $q:\mathcal S_{\rm cof}\to\mathcal S$
is fully faithful and surjective on objects.  Two lifts with the
same image have a unique comparison above that image's identity.
After placing the lifts on the same graph, that comparison is the
simultaneous gluing of the local three-boundary comparisons, by
full faithfulness of $q$.  The genus-zero comparison sends them
to identities, so (TC.13) descends through $q$.  Naturality makes
its component equal on any two such lifts.  This verifies both
existence and independence of descent as in
\cite[Propositions~6.7 and 6.12]{ET_BonattoRobertson2026}.
Thus the strict representative required for E.27 has one coherent
choice of comparisons, including all graph automorphisms.

The kernel calculation E.27 fixes the remaining one-holed-torus generator.
Its only completion inputs are the explicitly verified genus-one
normalizers E.26 and the Minamide--Nakamura classification
\cite[Theorem~4.3(i), Corollary~C]{ET_MN2022}, and the
last relation is the actual (G2) relation of the marked-pants operad.
The exact marked-pants relation is
\cite[Theorem~3.18, relation (G2), and equations (6.3)--(6.5)]
{ET_BonattoRobertson2026}.  After its inner ambiguity is written
$F(aba)=a^{-\eta}(aba)a^\eta$, that relation is
$1=d^{-\eta}a^\eta=(da^{-1})^{-\eta}$ in the free Birman kernel of
$B_4/\langle\omega_4\rangle$; abelianization forces $\eta=0$.
Consequently a kernel equivalence has one coherent identity comparison
on every recovered bordered object.  Normalize an arbitrary nodal
object, use these comparisons on its components, and clutch with the
fixed pins.  This defines one comparison depending only on the object.
Finite relabelling coherence descends it through its graph permutations.

The unique simultaneous section of $q,r$ makes all these comparisons
compatible with whole-arrow sewing.  In these coordinates the
four-point crossing is fixed; the preceding support construction fixes
every adjacent pinned crossing.  Endpoint equivariance fixes all its
translates.  The exact presentation (TC.8) then fixes every arrow and
composition.  This proves the injectivity claimed by E.42.
Existence is E.39 with (TC.3) and (TC.8) now supplied: on each adjacent
orbit use the common refining-quilt endpoint automorphisms, which
commute with both maps of its full stabilizer; on a diamond use one
quilt refining its finest multicurve.  The intermediate transports
cancel, so both composites coincide.  The same cap and sewing formulas
supply the operation comparisons.  The group law is the quilt cocycle.
Continuity follows at each finite categorical quotient and then by
inverse limit.  This completes the repair of the full-completion
comparisons needed in E.39--E.42, including genera three and higher;
no equality of $\widehat\Gamma$ and $\check\Gamma$ is used.
